Recent fault-characterization studies inject single-bit faults at gate
outputs and compare per-vector observability metrics across synthesis
flows. We show that, as commonly applied, these metrics can be
confounded by a variable that is not fault resistance: the granularity
of the gate-level representation. Using 15 cipher primitives from
PRESENT, GIFT, PRINCE, SIMON, Keccak-$\chi$, and Speck/ARX, we compare
truth-table LUT/AOIG, 2-input AIG, and MIG representations under the
same gate-output \bitflip{} model. The apparent benefit of majority
synthesis reverses sign with the baseline: relative to the LUT-style
AOIG baseline, MIG reduces average per-vector detection probability
(\AvgDP{}) by 13.0\% on average; relative to a decomposed-AIG baseline at comparable
small-gate granularity, MIG \emph{increases} \AvgDP{} by 21.4\%, with
no circuit showing a reduction. As a granularity control, we restrict
faults to output-driving gates (the realizable DFA target layer);
target-layer \AvgDP{} $=1.000$ across all three representations,
showing that the all-site difference comes from the interior fault-site
population, not from any change in the attacker-facing output layer. One finding remains
stable across representation: ARX primitives exhibit about $1.9\times$
broader output corruption per detected fault than substitution-based
primitives. We give reporting guidelines for synthesis-flow fault metrics
and release the netlists, campaign scripts, generated tables, and scripts
that check the quoted summary numbers.
